\documentclass[11pt,a4paper]{article}
\usepackage[a4paper,left=25mm,right=24mm,top=23mm,bottom=23mm,headheight=14pt,headsep=7mm,footskip=12mm]{geometry}
\usepackage{amsmath,amssymb,xcolor}
\usepackage[unicode,colorlinks=true,linkcolor=ink,urlcolor=ink]{hyperref}
\definecolor{ink}{rgb}{0.08,0.25,0.34}
\hypersetup{pdftitle={An Independent Verification of Null-Cone Quantum Gravity: Transition Amplitudes and Convergence},pdfauthor={Null-cone amplitude verification project},pdfsubject={A Reproducibility Study of Section 4.4 of arXiv:2611.34924v1}}
\ifdefined\pdfinfoomitdate\pdfinfoomitdate=1\fi
\ifdefined\pdftrailerid\pdftrailerid{}\fi
\ifdefined\pdfsuppressptexinfo\pdfsuppressptexinfo=15\fi
\setlength{\parindent}{0pt}
\setlength{\parskip}{5pt plus 0.4pt minus 0.4pt}
\setlength{\emergencystretch}{2em}
\raggedbottom
\makeatletter
\def\ps@audit{\def\@oddhead{\small\sffamily\color{ink}Null-cone amplitude verification\hfil Section 4.4}
\def\@oddfoot{\small\sffamily v0.1.0\hfil\thepage}\let\@evenhead\@oddhead\let\@evenfoot\@oddfoot}
\makeatother
\pagestyle{audit}
\newcommand{\topic}[1]{\section*{\color{ink}#1}}
\newcommand{\subtopic}[1]{\par\medskip{\large\bfseries #1}\par\nobreak}
\newcommand{\ket}[1]{\lvert #1\rangle}
\newcommand{\bra}[1]{\langle #1\rvert}
\newcommand{\norm}[1]{\left\lVert #1\right\rVert}
\newcommand{\Li}{\operatorname{Li}_2}
\newcommand{\E}{\mathbb E}
\newcommand{\paper}[2]{\href{https://arxiv.org/html/2611.34924v1\##1}{#2}}
\begin{document}

{\sffamily\small TECHNICAL REPORT\hfill 26 November 2026}\par
\vspace{4mm}
{\huge\bfseries\raggedright An Independent Verification of\\
Null-Cone Quantum Gravity:\\
Transition Amplitudes and Convergence\par}
\vspace{2mm}
{\large\itshape A Reproducibility Study of Section 4.4 of arXiv:2611.34924v1\par}
\vspace{4mm}

This study examines Eqs.~(4.29)--(4.33) of Venkatesa Chandrasekaran's v1 manuscript. For a single angular harmonic and polarization, the Gaussian state and mode rotation chosen in the paper give the connected pair amplitude
\[
 C(x)=\frac{x^2}{2}\Li(e^{-2ix}),\qquad
 x=\frac{1}{s}=\frac{4G_N\hbar}{\mathcal A}>0 .
\]
Here $\mathcal A$ is the fixed corner area. The sum over all radial frequencies converges absolutely for every fixed $x$. Its weak-coupling expansion contains $ix^3\log x$. Finiteness at positive coupling and the singular cubic coefficient follow from the same scalar function.

\textbf{Contributions and review.} OpenAI Codex performed the derivations and computations and drafted this report. This release has not been reviewed by an external specialist.

\subtopic{Claims and their dependencies}
\begin{tabular}{@{}p{0.20\linewidth}@{\hspace{0.035\linewidth}}p{0.72\linewidth}@{}}
\textbf{Paper equation} & \textbf{Result and hypotheses}\\[4pt]
(4.29), (4.30) & A complete orthonormal basis resolves one cut. Finite vector norms give absolute summability. Norm-preserving evolution gives the stated boundary-state bound.\\[4pt]
(4.31) & Contractive propagation and strong convergence of the cut projections control the error for any fixed finite number of cuts.\\[4pt]
(4.32) & The Gaussian widths, quadratic source and phase specify the preparation. The oscillator realization below implements those choices.\\[4pt]
(4.33) & Centering the source, orthogonality of pair states and the prescribed phase give the displayed sum. Absolute summability controls cutoff removal.
\end{tabular}

\subtopic{How to read the evidence}
\textbf{Proved} identifies an analytic implication with its hypotheses stated. \textbf{Conditional} marks a physical identification or preparation that supplies those hypotheses. \textbf{Numerical support} denotes finite calculations at specified precision. \textbf{Counterexample} identifies a universally phrased extension and a model that falsifies it. \textbf{Open interpretation} flags a specific unresolved identification. \textbf{Unverified} marks a question for which this audit supplies no proof.

Pages 2--5 give the derivations and bounds. Page 6 proves two counterexamples. Page 7 gives numerical evidence and page 8 records interpretive limits, provenance and review status. Source pagination refers to the \href{https://arxiv.org/pdf/2611.34924v1}{v1 PDF}, pp.~32--34.

\newpage
\topic{One cut and finitely many cuts}
\textbf{Proved.} Let $U_1:H_0\to H_\Sigma$ and $U_2:H_\Sigma\to H_f$ be bounded maps. Fix finite-norm vectors $\psi,\phi$ at the two ends. Set $a=U_1\psi$, $b=U_2^*\phi$, with inner products linear in the second argument. For a finite subset $F$ of a complete orthonormal basis on $H_\Sigma$, let $P_F$ be the orthogonal projection onto its span. Write $a_\alpha=\langle e_\alpha,a\rangle$ and likewise for $b$. Orthogonality gives
\[
 \langle b,P_Fa\rangle=\sum_{\alpha\in F}\overline{b_\alpha}a_\alpha,\qquad
 |\langle b,(I-P_F)a\rangle|
 \le \norm{(I-P_F)b}\,\norm{(I-P_F)a}.
\]
Completeness implies $P_Fa\to a$ and $P_Fb\to b$ in norm as $F$ increases. This proves the matrix identity in \paper{S4.E29}{(4.29)}. Applying Cauchy--Schwarz to the absolute values before taking the supremum over $F$ gives
\[
 \sum_\alpha |\overline{b_\alpha}a_\alpha|\le \norm b\,\norm a.
\]
For unitary $U_1,U_2$, the right side is $\norm\phi\,\norm\psi$, as in \paper{S4.E30}{(4.30)}. Contractions suffice for the inequality with these external norms.

\subtopic{The error with independent cutoffs}
Let $m$ be the fixed number of propagation segments. There are $m-1$ intermediate cuts. Suppose each $U_j:H_{j-1}\to H_j$ is a contraction. Define $a_0=b_0=\psi$, $a_j=U_j a_{j-1}$, and $b_j=P_j U_j b_{j-1}$ for $j<m$. On the final segment set $b_m=U_mb_{m-1}$. Subtract the recursions to get
\[
 a_j-b_j=(I-P_j)a_j+P_jU_j(a_{j-1}-b_{j-1}),\qquad j<m.
\]
Taking norms and iterating gives the estimate in \paper{S4.E31}{(4.31)},
\[
 |\langle\phi,a_m-b_m\rangle|
 \le \norm\phi\sum_{j=1}^{m-1}\norm{(I-P_j)U_j\cdots U_1\psi}.
\]
Each term on the right tests the projection on one \emph{fixed}, fully propagated vector. If each family of projections tends strongly to the identity, all finitely many terms tend to zero independently. This proves convergence with one cutoff per cut. It supplies no common bound as $m\to\infty$. The distinction from joint absolute summability at fixed $m$ already matters for two cuts (p.~6).

\textbf{Conditional.} Interpreting the product as the amplitude of a glued spacetime requires matching boundary data with opposite orientations, cancellation of the corresponding action terms, and the composition law for the admitted evolutions. The Hilbert-space argument starts once those maps are defined. The paper states these inputs immediately before (4.29).

\newpage
\topic{The Gaussian preparation and the pair}
\textbf{Conditional preparation.} Fix one real angular harmonic and one polarization. In the chosen product vacuum, the dimensionless radial coordinates are independent, $q_k\sim N(0,v_k)$, with zero mean momentum. A real wavefunction proportional to $\exp(-skq^2/2)$ has density proportional to $\exp(-skq^2)$, hence $v_k=(2sk)^{-1}=x/(2k)$. Take $Q_N=\sum_{k=1}^N q_k^2$. In the occupation basis set the vacuum phase to one and prescribe phase $e^{-ikx}$ per quantum. Additivity of the phase in occupation number then gives
\[
 U_N\ket{\boldsymbol n}
 =\exp\!\left[-ix\sum_{k\le N}kn_k\right]\ket{\boldsymbol n}.
\]
These three choices realize \paper{S4.E32}{(4.32)}. Each $U_N$ is unitary.

\subtopic{Independent coordinate derivation of (4.33)}
With $q=\sqrt{v_k}(a+a^\dagger)$ the normalized second oscillator wavefunction divided by its ground-state wavefunction is
$h_2(q)=(q^2/v_k-1)/\sqrt{2}$. Indeed the Gaussian moments
$\E q^2=v_k$, $\E q^4=3v_k^2$ give $\E h_2=0$, $\E h_2^2=1$. This is the even degree-two polynomial orthogonal to the vacuum. Consequently
\[
 (q_k^2-v_k)\ket0=\sqrt{2}\,v_k\ket{2_k},\qquad
 \bra0 U_N^*q_k^2U_N\ket{2_k}=\sqrt{2}\,v_k e^{-2ikx}.
\]
These are the creation and rotated annihilation matrix elements. For $B_N=Q_N-\sum_{k\le N}v_k$, set $b_N=B_N\ket0$. The vacuum is invariant. Subtracting the product of the two vacuum expectations makes the connected correlator $\langle b_N,U_N b_N\rangle$. Distinct $\ket{2_k}$ are orthogonal, so
\[
 C_N(x)=\sum_{k=1}^N 2v_k^2e^{-2ikx}
       =\frac{x^2}{2}\sum_{k=1}^N\frac{e^{-2ikx}}{k^2}.
\]
\textbf{Proved for this preparation.} $\norm{b_M-b_N}^2=2\sum_{N<k\le M}v_k^2\to0$. Thus $b_N\to b$ and $\norm b^2=\pi^2x^2/12$. Summing the absolutely convergent scalar series gives \paper{S4.E33}{(4.33)} with the closed-unit-disk definition of $\Li$.

\subtopic{Which unbounded operator is needed?}
In the vacuum sector the finite-occupation basis spans densely. Let $K$ act diagonally by $E_{\boldsymbol n}=\sum k n_k$ on the domain $\sum E_{\boldsymbol n}^2 |c_{\boldsymbol n}|^2<\infty$. This real diagonal operator is self-adjoint because testing its adjoint on each basis vector gives the same weighted square-summability condition. $U=e^{-ixK}$ is unitary everywhere. $U_N\to U$ strongly by eventual equality on each finite-occupation vector and uniform boundedness. The scalar limit is $\langle b,Ub\rangle$. Yet $\sum_k(2k)\,2v_k^2=x^2\sum_k1/k=\infty$, so $b\notin D(K^{1/2})$. Inserting the generator as a quadratic form on this vector requires a domain condition that fails.

\newpage
\topic{Absolute convergence and error bounds}
\textbf{Proved.} The absolute sum is $x^2\zeta(2)/2=\pi^2x^2/12$, including at $x=n\pi$. For a hard cutoff, integral comparison gives $\sum_{k>N}k^{-2}\le1/N$. If $\sin x\ne0$, the geometric sum of $z^k$, $z=e^{-2ix}$, on any consecutive interval has modulus at most $2/|1-z|=1/|\sin x|$. For $M=N+1$, partial summation against $k^{-2}$ bounds the infinite tail by
$(1/|\sin x|)\sum_{k=M}^\infty(k^{-2}-(k+1)^{-2})$. Thus
\[
 |C-C_N|\le\frac{x^2}{2}\min\!\left\{\frac1N,\frac{1}{|\sin x|(N+1)^2}\right\}.
\]
At exact resonance use the first bound. All terms there are positive and integral comparison from below gives $x^2/[2(N+1)]\le C-C_N\le x^2/(2N)$. Convergence is uniform on $0\le x\le X<\infty$ with the continuous value $C(0)=0$. It fails to be uniform on all $x>0$ by the unbounded sequence of exact resonances.

\subtopic{Hard, Abel and triangular cutoffs}
Let $H_N=\sum_{k\le N}1/k$. Any uniformly bounded weights tending to one mode by mode preserve the limit by domination by a constant times $k^{-2}$. For Abel weights $e^{-k/N}$, split at $N$ and use $1-e^{-k/N}\le k/N$ below the split. For triangular weights $(1-k/(N+1))_+$, the change below $N$ is exactly linear in $k$. The respective amplitude errors satisfy
\[
 |C-C_{\rm Abel}|\le\frac{x^2}{2N}(H_N+1),\qquad
 |C-C_{\rm tri}|\le\frac{x^2}{2}\left(\frac{H_N}{N+1}+\frac1N\right).
\]
The cubic coefficients divided by $-i$ before cutoff removal are $H_N$, $-\log(1-e^{-1/N})$ and $H_N-N/(N+1)$ for hard, Abel and triangular weights. This follows by weighting $1/k$ and summing, using $\sum r^k/k=-\log(1-r)$. After subtracting $\log N$ the constants tend to $\gamma,0,\gamma-1$, respectively. The common amplitude and the cutoff-dependent perturbative coefficients answer different questions.

\subtopic{A second expression for the exact hard-cutoff error}
Write $1/k^2=\int_0^\infty u e^{-ku}\,du$. Absolute integrability after summing follows from $\sum_{k>N}1/k^2<\infty$. Summing the geometric series inside the integral and setting $y=Mu$, $M=N+1$, gives
\[
 C-C_N=\frac{x^2z^M}{2M^2}\int_0^\infty
              \frac{y e^{-y}}{1-z e^{-y/M}}\,dy .
\]
The integrand tends to $M$ at the lower endpoint when $z=1$ and to $0$ otherwise. This formula permits a numerical check using only a finite sum and quadrature.

\newpage
\topic{What fails in the weak-coupling expansion?}
Keep $\mathcal A,\hbar>0$ fixed and vary $x=4\hbar G_N/\mathcal A$. For each finite $N$, Taylor expansion gives $C_N=x^2 H_N^{(2)}/2-ix^3H_N+O_N(x^4)$, where $H_N^{(2)}=\sum_{k\le N}k^{-2}$. The cubic coefficient has no finite limit. To expand the summed function, put $F(t)=\sum_{k\ge1}e^{-ikt}/k^2$. With Abel factor $r^k$, $0<r<1$, differentiating gives $F_r'(t)=i\log(1-r e^{-it})$. As $r\uparrow1$ the derivative converges uniformly on compact subintervals of $(0,2\pi)$ and the functions converge uniformly on the whole circle. The principal argument of $1-e^{-it}$ is $(\pi-t)/2$ there. Integration from a positive lower endpoint, then to zero using the integrable $\log t$ singularity, yields
\[
 F(t)=\frac{\pi^2}{6}-\frac{\pi t}{2}+\frac{t^2}{4}
      +i\int_0^t\log(2\sin(u/2))\,du .
\]
For $0<x<\pi$ take $t=2x$. The Euler product for the sine yields the uniformly convergent series on $0\le y\le x$,
\[
 \log(\sin y/y)=-\sum_{j\ge1}\zeta(2j)(y/\pi)^{2j}/j .
\]
Integrating term by term gives
\begin{align*}
 \Re C(x)&=\pi^2x^2/12-\pi x^3/2+x^4/2,\\
 \Im C(x)&=x^3\left[\log(2x)-1-\sum_{j\ge1}
       \frac{\zeta(2j)(x/\pi)^{2j}}{j(2j+1)}\right].
\end{align*}
For $J\ge0$, bound the positive tail by replacing $\zeta(2j)$ by $\zeta(2)$, $j(2j+1)$ by its first remaining value, and summing the geometric powers. The absolute error in $\Im C$ is at most
\[
 \frac{x^3\zeta(2)(x/\pi)^{2J+2}}
 {(J+1)(2J+3)(1-(x/\pi)^2)}.
\]
The first two corrections to $x^3[\log(2x)-1]$ are $-x^5/18-x^7/900$. The regular power series can be differentiated inside its radius of convergence. Thus $C'''(x)=-3\pi+i[6\log(2x)+5]+O(x)$ at $0^+$. The third derivative in $G_N$ has the extra factor $(4\hbar/\mathcal A)^3$.

\subtopic{Values versus coefficients}
\textbf{Proved.} The bound $|C_N|\le\pi^2 x^2/12$ gives zero in both orders for the value as $N\to\infty$, $x\downarrow0$. Define
$T_N=[\Re C_N-x^2 H_N^{(2)}/2]/x^3$. At fixed $N$, the finite cosine sum gives $T_N\to0$. At fixed $x>0$, $T_N\to-\pi/2+x/2$ for $x<\pi$. The two orders for this cubic coefficient therefore give $0$ and $-\pi/2$. The divergent imaginary coefficient and the order-dependent real coefficient locate the failure precisely. This differentiates a family of prepared \emph{scalar} amplitudes. Comparing states or unbounded operators across couplings needs an identification of their representations.

\newpage
\topic{Two counterexamples and their scope}
\subtopic{1. The uncentered positive source in the same vacuum}
\textbf{Claim tested.} The positive forms of $Q_N$ from (4.32) have a finite pointwise limit on some dense domain in the product Gaussian representation.

\textbf{Counterexample.} Use exactly $v_k=x/(2k)$ at any fixed $x>0$ and let $\mu$ be the product measure. Integrating a Gaussian against $e^{-q_k^2/x}$ gives $(1+2v_k/x)^{-1/2}$. Independence now makes the many-mode transform telescope,
\[
 \E_\mu e^{-Q_N/x}=\prod_{k=1}^N \sqrt{k/(k+1)}=(N+1)^{-1/2}.
\]
The random variables $Q_N$ increase to an extended nonnegative $Q_\infty$. By dominated convergence $\E e^{-Q_\infty/x}=0$. The integrand is nonnegative, so $Q_\infty=\infty$ almost surely. For each nonzero $f\in L^2(\mu)$ the set $|f|>0$ has positive measure. Monotone convergence gives $\lim_N\int Q_N|f|^2\,d\mu=\infty$. Hence the finite-form domain is exactly $\{0\}$.

\textbf{Scope.} The connected prescription has an $L^2$ limit because the variance of $B_M-B_N$ is $2\sum_{N<k\le M}v_k^2\to0$. Its real limit $B$ defines a multiplication operator on $\{f:Bf\in L^2(\mu)\}$, which contains the vacuum $f=1$. This construction explains how (4.33) survives the failure of the indicated raw positive limit.

\subtopic{2. Joint absolute summability with two cuts}
\textbf{Claim tested.} Finite-norm external states and unitary evolution suffice for absolute summability of all paths after expanding two intermediate cuts in orthonormal bases.

\textbf{Counterexample.} Take $H=\bigoplus_{j\ge1}\mathbb C^{d_j}$, $d_j=4^j$, with standard coordinate bases and normalized constant block vectors $u_j$. Set $\psi=\phi=\sum_{j\ge1}2^{-j/2}u_j$, so both norms are one. Use the three maps $I,W,I$. On block $j$, $W_{ab}=(-1)^{a\cdot b}/\sqrt{d_j}$ for binary strings of length $2j$. Orthogonality follows by pairing opposite signs in the sum over strings whenever two rows differ. Thus the direct sum $W$ is unitary.

Every within-block path has modulus $2^{-j}/d_j^{3/2}$. There are $d_j^2$ such paths. Their absolute mass is $2^{-j}\sqrt{d_j}=1$ in each block, so total absolute mass diverges. On the other hand $Wu_j$ is the first coordinate vector. Its overlap with $u_j$ is $1/\sqrt{d_j}$. The actual matrix element is $\sum_{j\ge1}2^{-j}/\sqrt{d_j}=\sum_{j\ge1}4^{-j}=1/3$.

\textbf{Scope.} This Hilbert-space example satisfies both the one-cut theorem and the finite-cut projection estimate. It disproves the stated universal extension to unrestricted rearrangements of individual paths. Neither example decides ultraviolet completeness in four-dimensional gravity.

\newpage
\topic{Numerical checks and a simultaneous limit}
\textbf{Numerical support.} The release contains 11 unit tests and 127 numerical records. Every tested input and the printed complex outputs are in \texttt{baseline/numerical.jsonl}. The hard-cutoff evaluations use a direct complex sum. The reference evaluations use the dilogarithm, the real log-sine integral on p.~5 and, at five inputs, a direct prefix plus the Laplace tail on p.~4. The last two methods do not call the dilogarithm. They share the mpmath arithmetic library.

\begin{center}
\begin{tabular}{@{}r@{\qquad}r@{\qquad}r@{}}
$s$ & $\Re C(1/s)$ & $\Im C(1/s)$\\ \hline
1   & $-0.2483292933707834$ & $-0.3635730254316396$\\
10  & $\phantom{-}0.006703874007446236$ & $-0.002609993579151187$\\
100 & $\phantom{-}0.00008068090701561643$ & $-0.000004912028560994813$
\end{tabular}
\end{center}
At $s=10$, $N=2048$, the observed absolute error is $5.96735\,10^{-9}$ and the analytical upper bound is $1.19292\,10^{-8}$ (rounded displays). At 125 decimal digits, the largest observed discrepancy divided by $x^2$ was $1.245\,10^{-120}$ or less for nine log-sine comparisons and $3.613\,10^{-124}$ or less for the five Laplace comparisons. These discrepancies measure agreement between calculations. They are not rigorous roundoff enclosures.

The nine inputs include $x=10^{-18}$, $\pi/2$, $\pi$, $\pi\pm10^{-25}$ and $10^6+1/7$. Symbolic Gaussian integration and a separate number-basis multiplication check the sign and both factors of $\sqrt{2}$. A standard-library script also checks $2C(\pi/4)/(\pi/4)^2=-\pi^2/48-i\,\mathrm{Catalan}$ by grouping phases modulo four and evaluating the alternating odd reciprocal-square sum. It enumerates the first four Walsh blocks and checks the Gaussian product using exact rational arithmetic.

\subtopic{When the cutoff probes the oscillation scale}
\textbf{Proved.} Write $x=\lambda/N$, $\lambda>0$. Removing the quadratic and finite-cutoff cubic terms leaves
\[
 \frac{C_N-x^2H_N^{(2)}/2}{x^3}+iH_N
   =x\sum_{k=1}^N h(kx),\qquad
 h(u)=\frac{e^{-2iu}-1+2iu}{2u^2}
     =-2\int_0^1(1-t)e^{-2iut}\,dt .
\]
The integral form is Taylor's formula with integral remainder. It gives $h(0)=-1$ and $|h'|\le4\int_0^1t(1-t)dt=2/3$. On each interval of length $x$, the right-endpoint quadrature error is bounded by $x^2/3$. Summing gives $\lambda^2/(3N)$ for the difference from $\int_0^\lambda h(u)\,du$.

\textbf{Numerical support.} A separately coded rational Taylor integral checks $\lambda=1/4,1,8$ at $N=64,256,1024,4096$. At $\lambda=1$, $N=4096$ the observed error is $7.55088\,10^{-5}$, below $8.13802\,10^{-5}$. The scale $kx=O(1)$ explains why expanding every mode at fixed cutoff can miss a term in the summed amplitude.

\newpage
\topic{Physical interpretation and provenance}
\textbf{Conditional.} The paper uses causal factorization and unitary evolution for the admitted deformations to identify the glued amplitude. It allows positive Gaussian widths in its representation and selects the explicit widths and phases in (4.32). Appendix D.3 motivates a short-wavelength scaling with fixed endpoint data. Scaling alone does not fix the coefficient at every $k$. \textbf{Open interpretation.} Calling the prescribed number rotation an actual translation on a finite null boundary would require a geometric map preserving its endpoint data. Equation (4.33) needs only the prescribed rotation.

\textbf{Unverified.} Appendix C.2 gives sufficient conditions for bounded self-adjoint, strongly continuous operator families with integrable operator norm. The restricted calculation supplies no continuum self-adjointness or common-domain construction for every unbounded gravitational source. The cut proof has no logical cycle as a Hilbert-space implication. Using that estimate by itself to justify the existence or unitarity of the same physical evolution would assume a premise in the conclusion.

The example sums radial frequencies without bound for one prepared, centered observable. A claim of ultraviolet completeness would need control of a specified class of physical observables, their domains, removal of regulators and consistent dynamics at arbitrarily short distances. These additional propositions are \textbf{unverified} here. The calculation gives no verdict on the full theory.

\subtopic{Contributions and review}
The project lead defined the question and scope and authorized publication. OpenAI Codex performed the primary analytic derivations, designed and ran the symbolic and numerical checks, and prepared this report and software. During the research stage a second AI context inspected the primary manuscript and recorded its reasoning before reading the first analysis. It tested the raw source and the two-cut extension and wrote the separate standard-library checks. Shared mathematical assumptions and potential model errors remain possible. The evidence for each theorem is its displayed proof.

As of 26 November 2026, this verification has not been reviewed by an external specialist. The project record contains no attestation of a human line-by-line mathematical audit. Venkatesa Chandrasekaran authored the manuscript under study. Project title and version identify this independent work.

\subtopic{Primary sources and version}
V. Chandrasekaran, \textit{An algebraic quantization of gravity from the null cone}, arXiv:2611.34924v1, 25 November 2026. The five numbered equations appear on printed pp.~32--34. Context for gluing and source selection is linked above. \href{https://arxiv.org/abs/2611.34924v1}{Version record} $\cdot$ \paper{S4.SS4}{Section 4.4} $\cdot$ \paper{A3.SS2}{Appendix C.2} $\cdot$ \paper{A4.SS3}{Appendix D.3}.

NIST Digital Library of Mathematical Functions, \href{https://dlmf.nist.gov/25.12.E1}{Eq.~25.12.1} (dilogarithm series) and \href{https://dlmf.nist.gov/4.22.E1}{Eq.~4.22.1} (sine product). These identities fix the boundary value and the convergent kernel expansion used on pp.~3--5.

\texttt{reference/paper.json} records both arXiv files' lengths and SHA-256 digests. \texttt{SHA256SUMS} covers the report source, PDF, code and reference results.

\end{document}
